\documentclass[../../main.tex]{subfiles}
\graphicspath{{\subfix{../../image/}}}

\begin{document}

\section{Bernoulli 试验场合的极限定理}\label{section:ch05-Bernoulli试验场合极限定理}

\subsection*{一、问题的提出}

在第一章中我们已经指出：人们在长期实践中发现，虽然个别随机事件在某次试验中可以出现也可以不出现，但是在大量重复试验中却呈现出明显的规律性，即一个随机事件出现的频率在某个固定数的附近摆动，这就是所谓``频率稳定性''。对于这点，迄今为止，我们尚未给予理论上的说明。

数学上怎样来描述在一定条件下的大量重复试验呢？我们在第二章已经建立了 Bernoulli 试验这一概率模型，并指出它可以作为在一定条件下的重复试验的数学模型。在 Bernoulli 试验中，各次试验是相互独立的，并且在每次试验中，我们所关心的事件 $A$ 出现的概率 $P(A)=p$ 保持不变，这些特征可以看作是从数学角度把``在一定条件下''、``重复试验''等等用语的含义加以明确化。

在 Bernoulli 试验中，若以 $\mu_n$ 记 $n$ 次试验中 $A$ 出现的次数，则 $\dfrac{\mu_n}{n}$ 便是在这 $n$ 次试验中事件 $A$ 出现的频率，所谓频率稳定性无非是指当试验次数 $n$ 增大时，频率 $\dfrac{\mu_n}{n}$ 接近于某个固定的常数。

这个固定的常数就是事件 $A$ 在一次试验中发生的概率。由此可见，讨论频率 $\dfrac{\mu_n}{n}$ 的极限行为是理解概率论中最基本的概念——概率所不可缺少的。正是这个缘故，在概率论的发展史上，极限定理的研究一直占重要地位，而它的发源地就是 Bernoulli 试验这个概型。

从前几章的讨论中我们知道，$\mu_n$ 是随机变量，它服从二项分布
\begin{align}
P\{\mu_n=k\}=\binom{n}{k}p^{k}q^{n-k},\qquad k=0,1,2,\cdots,n
\end{align}
其数学期望 $\boldsymbol{E}\mu_n=np$，方差 $\boldsymbol{D}\mu_n=npq$。这在一定程度上帮助我们进一步了解了频率 $\dfrac{\mu_n}{n}$ 的性质。但是我们更需要知道的是 $n$ 很大时 $\mu_n$ 或 $\dfrac{\mu_n}{n}$ 的性质。

显然，当 $n$ 很大时，$\mu_n$ 一般也很大，所以直接研究 $\mu_n$ 不很恰当，还是研究频率 $\dfrac{\mu_n}{n}$ 为宜。因为 $\boldsymbol{E}\!\left(\dfrac{\mu_n}{n}\right)=p$，$\boldsymbol{D}\!\left(\dfrac{\mu_n}{n}\right)=\dfrac{pq}{n}$，所以当 $n\to\infty$ 时，频率的数学期望保持不变，而方差则趋于 $0$。我们知道方差为 $0$ 的随机变量是常数，于是我们自然预期频率将趋于常数 $p$（即事件 $A$ 发生的概率）。但是频率 $\dfrac{\mu_n}{n}$ 是随机变量，关于它的极限又将采用何种提法呢？

一种提法是：当 $n$ 足够大时，频率 $\dfrac{\mu_n}{n}$ 与概率 $p$ 有较大偏差的概率很小。用数学语言来讲，就是要证明：对于任意 $\varepsilon>0$，
\begin{align}
\lim_{n\to\infty}P\!\left\{\left|\frac{\mu_n}{n}-p\right|\geqslant\varepsilon\right\}=0\label{align:ch05-频率大偏差趋于零}
\end{align}
或者它的等价的式子成立，即
\begin{align}
\lim_{n\to\infty}P\!\left\{\left|\frac{\mu_n}{n}-p\right|<\varepsilon\right\}=1
\end{align}

历史上，雅·Bernoulli 第一个研究了这种类型的极限定理，在他死后于 1713 年发表的论文中（这是概率论的第一篇论文！）建立了\cref{align:ch05-频率大偏差趋于零}式，这是一大类概率论极限定理——大数定律（law of large numbers）中的第一个。

关于频率接近于概率还有其他提法，譬如近二百年后 Borel 建立了
\begin{align}
P\!\left\{\lim_{n\to\infty}\frac{\mu_n}{n}=p\right\}=1
\end{align}
从而开创了另一种形式的极限定理——强大数定律的研究。

本节只讨论 Bernoulli 大数定律，关于强大数定律的讨论将在 \S4 中进行。

为了研究 $\mu_n$ 的极限行为，可以讨论它的分布 $P\{\mu_n<x\}$ 的变化情况。但是由于 $\boldsymbol{E}\mu_n=np$，$\boldsymbol{D}\mu_n=npq$，因此对于固定的 $x$ 来考虑 $P\{\mu_n<x\}$ 的极限不会有多大意义，因为它将趋于 $0$，所以通常改为研究``标准化''的随机变量
\begin{align}
\zeta_n=\frac{\mu_n-np}{\sqrt{npq}}\label{align:ch05-标准化随机变量}
\end{align}
的分布函数
\begin{align*}
P\{\zeta_n<x\}
\end{align*}
的极限行为，由 $\zeta_n$ 的分布函数不难求得 $\mu_n$ 的分布函数。

关于上述分布，已证明它的极限分布是正态分布 $N(0,1)$，即
\begin{align}
\lim_{n\to\infty}P\{\zeta_n<x\}=\frac{1}{\sqrt{2\pi}}\int_{-\infty}^{x}\mathrm{e}^{-t^{2}/2}\,\mathrm{d}t
\end{align}
这个结果最早由法国数学家 De Moivre（1667—1754）于 1718 年建立，他对 $p=\dfrac12$ 证得了上述结果。后来，1812 年由 Laplace 推广到 $0<p<1$ 的一般场合，那是另一类概率论极限定理——中心极限定理中的第一个。\footnote{关于 De Moivre 发现正态逼近的年份众说纷纭，这里采用费勒的说法，见 [2]174 页。}

下面将会看到极限定理的研究直接联系到大 $n$ 场合的二项分布的计算，这就顺便解决了第二章遗留下来的一些计算问题。

形如上述收敛于正态分布的极限定理的研究，在长达两个世纪的时期内成了概率论研究的中心课题，因此在上世纪 20 年代由波利亚命名为中心极限定理。

经过长期的研究，人们认识到，$\mu_n$ 具有上述性质是由于它是独立随机变量之和。事实上，若令
\begin{align}
\xi_i=
\begin{cases}
1,&\text{第 }i\text{ 次试验出现 }A,\\
0,&\text{第 }i\text{ 次试验不出现 }A,
\end{cases}\label{align:ch05-伯努利指示变量}
\end{align}
则
\begin{align}
\mu_n=\xi_1+\xi_2+\cdots+\xi_n
\end{align}
这里 $\xi_1,\xi_2,\cdots,\xi_n$ 是相互独立的。以后我们将会看到，对于一般的随机变量 $\xi_i$，$i=1,2,\cdots,n$，也可以研究它们的和的极限定理，并且在一定条件下，这个和也具有类似于 $\mu_n$ 的性质，关于这些问题的研究就构成了本章的主要内容。

为了叙述方便起见，我们引进如下定义。

\begin{definition}[(弱)大数定律]\label{definition:ch05-弱大数定律}
称随机变量序列 $\{X_n,n\geqslant1\}$ 服从（弱）大数定律，若存在数列 $\{a_n\}$ 和 $\{b_n\}$，$b_n\to+\infty$，使
\begin{align*}
\frac{S_n-a_n}{b_n}\xrightarrow{P}0,
\end{align*}
其中 $S_n=\sum_{k=1}^{n}X_k$。
\end{definition}

\begin{note}
通常不作特别说明时，指 $a_n=\boldsymbol{E}S_n$，$b_n=n$ 的规范形式。
\end{note}

\begin{definition}[(弱)大数定律的规范形式]\label{definition:ch05-弱大数定律规范形式}
称 $\{X_n\}$ 服从（弱）大数定律，若
\begin{align*}
\frac{S_n}{n}-\frac{\boldsymbol{E}S_n}{n}\xrightarrow{P}0,
\end{align*}
即
\begin{align*}
\lim_{n\to\infty}P\!\left\{\left|\frac1n\sum_{k=1}^{n}X_k-\frac1n\sum_{k=1}^{n}\boldsymbol{E}X_k\right|<\varepsilon\right\}=1.
\end{align*}
\end{definition}

\begin{definition}[大数定律]\label{definition:ch05-大数定律}
若 $\xi_1,\xi_2,\cdots,\xi_n,\cdots$ 是随机变量序列，令
\begin{align}
\eta_n=\frac{\xi_1+\xi_2+\cdots+\xi_n}{n}\label{align:ch05-算术平均}
\end{align}
如果存在这样的一个常数序列 $a_1,a_2,\cdots,a_n,\cdots$，对任意的 $\varepsilon>0$，恒有
\begin{align}
\lim_{n\to\infty}P\{|\eta_n-a_n|<\varepsilon\}=1
\end{align}
则称序列 $\{\xi_n\}$ 服从大数定律（或大数法则）。
\end{definition}

在以后的讨论中，我们几乎总是假定 $\xi_1,\xi_2,\cdots,\xi_n,\cdots$ 是独立随机变量序列，显然，Bernoulli 大数定律是一般大数定律的一种特殊场合。

关于中心极限定理，我们总是对独立随机变量序列 $\xi_1,\xi_2,\cdots,\xi_n,\cdots$ 进行讨论，假定 $\boldsymbol{E}\xi_i$ 及 $\boldsymbol{D}\xi_i$ 存在，令
\begin{align}
\zeta_n=\frac{\sum_{i=1}^{n}\xi_i-\sum_{i=1}^{n}\boldsymbol{E}\xi_i}{\sqrt{\sum_{i=1}^{n}\boldsymbol{D}\xi_i}}\label{align:ch05-标准化和}
\end{align}
我们的目的是寻找使
\begin{align}
\lim_{n\to\infty}P\{\zeta_n<x\}=\frac{1}{\sqrt{2\pi}}\int_{-\infty}^{x}\mathrm{e}^{-t^{2}/2}\,\mathrm{d}t
\end{align}
成立的条件。

一般，若独立随机变量序列 $\xi_1,\xi_2,\cdots,\xi_n,\cdots$ 的标准化和 $\zeta_n$ 使上式成立，则我们称 $\{\xi_i\}$ 服从中心极限定理（central limit theorem）。

显然，De Moivre--Laplace 极限定理是中心极限定理的特例。

\subsection*{二、Bernoulli 大数定律}

我们将证明一个比 Bernoulli 大数定律更强的命题。

\subsubsection*{1.\ Chebyshev 大数定律}

\begin{theorem}[Chebyshev 大数定律]\label{theorem:ch05-Chebyshev大数定律}
设 $\xi_1,\xi_2,\cdots,\xi_n,\cdots$ 是由两两不相关的随机变量所构成的序列，每一随机变量都有有限的方差，并且它们有公共上界
\begin{align*}
\boldsymbol{D}\xi_1\leqslant C,\quad \boldsymbol{D}\xi_2\leqslant C,\quad\cdots,\quad \boldsymbol{D}\xi_n\leqslant C,\quad\cdots
\end{align*}
则对任意的 $\varepsilon>0$，皆有
\begin{align}
\lim_{n\to\infty}P\!\left\{\left|\frac1n\sum_{k=1}^{n}\xi_k-\frac1n\sum_{k=1}^{n}\boldsymbol{E}\xi_k\right|<\varepsilon\right\}=1\label{align:ch05-Chebyshev大数定律结论}
\end{align}
\end{theorem}

\begin{proof}
因为 $\{\xi_k\}$ 两两不相关，故
\begin{align*}
\boldsymbol{D}\!\left(\frac1n\sum_{k=1}^{n}\xi_k\right)=\frac{1}{n^{2}}\sum_{k=1}^{n}\boldsymbol{D}\xi_k\leqslant\frac{C}{n}.
\end{align*}
再由 Chebyshev 不等式得到
\begin{align*}
P\!\left\{\left|\frac1n\sum_{k=1}^{n}\xi_k-\frac1n\sum_{k=1}^{n}\boldsymbol{E}\xi_k\right|<\varepsilon\right\}
\geqslant 1-\frac{\boldsymbol{D}\!\left(\frac1n\sum_{k=1}^{n}\xi_k\right)}{\varepsilon^{2}}
\geqslant 1-\frac{C}{n\varepsilon^{2}}.
\end{align*}
所以
\begin{align*}
1\geqslant P\!\left\{\left|\frac1n\sum_{k=1}^{n}\xi_k-\frac1n\sum_{k=1}^{n}\boldsymbol{E}\xi_k\right|<\varepsilon\right\}\geqslant 1-\frac{C}{n\varepsilon^{2}}.
\end{align*}
于是，当 $n\to\infty$ 时有\cref{align:ch05-Chebyshev大数定律结论}，因此定理得证。
\end{proof}

这个结果在 1866 年被俄国数学家 Chebyshev 所证明，它是关于大数定律的一个相当普遍的结论，许多大数定律的古典结果是它的特例；此外，证明这个定律所用的方法后来称为矩法，也很有创造性，在这基础上发展起来的一系列不等式是研究各种极限定理的有力工具。

Markov（1856—1922）注意到在 Chebyshev 的论证中，只要
\begin{align}
\frac{1}{n^{2}}\boldsymbol{D}\!\left(\sum_{k=1}^{n}\xi_k\right)\to 0\label{align:ch05-Markov条件}
\end{align}
则大数定律就能成立，通常称条件\cref{align:ch05-Markov条件}为 Markov 条件。

\begin{theorem}[Markov 大数定律]\label{theorem:ch05-Markov大数定律}
对于随机变量序列 $\xi_1,\xi_2,\cdots,\xi_n,\cdots$，若\cref{align:ch05-Markov条件}成立，则对任意 $\varepsilon>0$，均有\cref{align:ch05-Chebyshev大数定律结论}。
\end{theorem}

Chebyshev 大数定律显然可由 Markov 大数定律推出；更重要的是 Markov 大数定律已经没有任何关于独立性的假定。研究相依随机变量序列的大数定律是近代概率论的课题之一，但是这已超出我们讨论的范围。

\subsubsection*{2.\ Bernoulli 大数定律与 Poisson 大数定律}

\begin{theorem}[Bernoulli 大数定律]\label{theorem:ch05-Bernoulli大数定律}
设 $\mu_n$ 是 $n$ 次 Bernoulli 试验中事件 $A$ 出现的次数，而 $p$ 是事件 $A$ 在每次试验中出现的概率，则对任意 $\varepsilon>0$，都有
\begin{align}
\lim_{n\to\infty}P\!\left\{\left|\frac{\mu_n}{n}-p\right|<\varepsilon\right\}=1
\end{align}
\end{theorem}

\begin{proof}
定义随机变量 $\xi_i$ 如\cref{align:ch05-伯努利指示变量}，则
\begin{align*}
\boldsymbol{E}\xi_i=p,\qquad \boldsymbol{D}\xi_i=pq\leqslant\frac14,
\end{align*}
而
\begin{align*}
\frac1n\sum_{k=1}^{n}\xi_k-\frac1n\sum_{k=1}^{n}\boldsymbol{E}\xi_k=\frac{\mu_n}{n}-p,
\end{align*}
故由 Chebyshev 大数定律立刻推出 Bernoulli 大数定律。
\end{proof}

显然，Bernoulli 大数定律也可以通过 Chebyshev 不等式直接加以证明：
\begin{align}
P\!\left\{\left|\frac{\mu_n}{n}-p\right|\geqslant\varepsilon\right\}
\leqslant\frac{1}{\varepsilon^{2}}\boldsymbol{D}\!\left\{\frac{\mu_n}{n}\right\}
=\frac{1}{n\varepsilon^{2}}\boldsymbol{D}\xi_i
\leqslant\frac{1}{4n\varepsilon^{2}}
\end{align}

历史上，Bernoulli 是通过直接展开和繁复的计算才证得大数定律的。

\begin{theorem}[Poisson 大数定律]\label{theorem:ch05-Poisson大数定律}
如果在一个独立试验序列中，事件 $A$ 在第 $k$ 次试验中出现的概率等于 $p_k$，以 $\mu_n$ 记在前 $n$ 次试验中事件 $A$ 出现的次数，则对任意 $\varepsilon>0$，都有
\begin{align}
\lim_{n\to\infty}P\!\left\{\left|\frac{\mu_n}{n}-\frac{p_1+p_2+\cdots+p_n}{n}\right|<\varepsilon\right\}=1
\end{align}
\end{theorem}

\begin{proof}
定义 $\xi_k$ 为第 $k$ 次试验中事件 $A$ 出现的次数，则
\begin{align*}
\boldsymbol{E}\xi_k=p_k,\qquad \boldsymbol{D}\xi_k=p_k(1-p_k)\leqslant\frac14,
\end{align*}
再用 Chebyshev 大数定律立刻可以推出结论。
\end{proof}

Poisson 提出了不同于 Bernoulli 试验的另一种独立试验模型，证明了二项分布的逼近定理，导出 Poisson 分布（皆见第二章 \S4）以及证明了上述大数定律，从而奠定了他在概率论发展史上的重要地位。顺便指出，也是他第一个使用``大数定律''这一名称。

纵观大数定律发展史，对科学发展规律的理解，启发良多。首要是新概念的提出，从 Bernoulli 的首创，到 Poisson 的推广，再到 Chebyshev，模型越来越普遍。其次是方法和数学工具的进步，从 Bernoulli 的直接估算，到 Chebyshev 的矩法，显然后者有力得多。还有，细致的考察也十分重要。Markov 通过对老师 Chebyshev 证明细节的深入研究，提出使用最普遍而又摆脱独立性的条件，成集大成者。这些规律为后来的发展一再证实。

\subsubsection*{3.\ 大数定律的重要意义}

Bernoulli 大数定律建立了在大量重复独立试验中事件出现频率的稳定性，正因为这种稳定性，概率的概念才有客观意义。

观察个别现象时是连同一切个别的特性来观察的。这些个别的特性往往蒙蔽了事物的规律性。通过平均，在大量观察中个别因素的影响将相互抵消而使总体稳定。例如，虽然每个气体分子的运动带有很大的随机性，但是作为气体平均特征的压力、温度等却是稳定的，大数定律说明了这种稳定性。古人用``定律''来称呼这类命题，大概认为它与物理学中的运动三大定律，万有引力定律，化学中的定比定律等一样是宇宙固有的规律。

从现代的观点来看，大数定律也是一类数学定理，它有一定的条件和假设，其结论可以通过通常的数学方法证明。它的要点在于：一、要求 $n$ 很大，即它是一类极限定理。在数学中，极限定理有的是；二、关于平均值，这就有点概率论特色；三、建立概率接近于 $1$ 或 $0$ 的规律，这是概率论研究中特别强调的；四、规律的产生是大量独立或弱相关因素积累的结果，这就涉及统计独立性这一概率论特有的概念。因此大数定律是概率论这一学科中最有特色的命题。

总之，大数定律在偶然性与必然性之间架起了桥梁，对人类认识客观世界大有启迪，是自然哲学的重要组成部分。另一方面，大数定律也有许多应用。例如，Bernoulli 大数定律还提供了通过试验来确定事件概率的方法，既然频率 $\dfrac{\mu_n}{n}$ 与概率 $p$ 有较大偏差的可能性很小，那么我们便可以通过做试验确定某事件发生的频率并把它作为相应概率的估计。这类方法称为参数估计，它是数理统计中的主要研究课题之一，参数估计的重要理论基础之一就是大数定律。

下面我们以保险业为例来说明人类如何利用大数定律以增进社会福祉。

天有不测风云，人有旦夕祸福。自然界和社会生活都充满着不确定性。一场灾害或事故可使一家人立即陷入经济困境，个人对此无能为力，但社会却可建立损失分担机制，保险业与精算学应运而生。作为精算科学基础的保险原理，即损失分担原理，可以认为是大数定律用保险术语的重述。

精算学的研究主题是分析巨大的、不可预测的损失的各种金融后果，并设计某种机制以缓冲这类损失的有害的金融效应。例如在财产保险中每月支付一笔小钱以求在一场大火或失窃后得到高额赔偿，便是人类对抗偶然性的有效方法。

为规避某类风险，人们找中介即保险公司，保险公司则创造满足下面 $4$ 个条件的群体：
\begin{enumerate}
\item 损失是不可预测的；
\item 风险是独立的；
\item 风险是齐性的；
\item 这个群体相当大，使得各个个体要求赔偿的整个损失额变成相对确定。
\end{enumerate}

读者不难把上述 $4$ 个条件与前面讲过的或后面将要讲到的大数定律成立的条件作对比。

因此可以说，保险业是人类利用大数定律的范例。

\subsection*{三、De Moivre--Laplace 极限定理}

大数定律只断言 $P\!\left\{\left|\dfrac{\mu_n}{n}-p\right|\geqslant\varepsilon\right\}$ 当 $n\to\infty$ 时趋于 $0$，也即 $\dfrac{\mu_n}{n}$ 接近于 $p$，而 De Moivre--Laplace 极限定理则给出 $\mu_n$ 的渐近分布的更精确表述。

下面定理给出了两个结果。第一个结果提供了 $P\{\mu_n=k\}$ 的渐近表达式，这类结果一般称为局部极限定理。第二个结果给出了标准化随机变量 $\dfrac{\mu_n-np}{\sqrt{npq}}$ 的渐近分布，称为积分极限定理，它是一般中心极限定理的特例。这两个结果既有区别也有联系。

\begin{theorem}[De Moivre--Laplace]\label{theorem:ch05-DeMoivreLaplace}
若 $\mu_n$ 是 $n$ 次 Bernoulli 试验中事件 $A$ 出现的次数，$0<p<1$，则对任意有限区间 $[a,b]$：
\begin{enumerate}
\item 当 $a\leqslant x_k\equiv\dfrac{k-np}{\sqrt{npq}}\leqslant b$ 及 $n\to\infty$ 时，一致地有
\begin{align}
P\{\mu_n=k\}\div\left(\frac{1}{\sqrt{npq}}\cdot\frac{1}{\sqrt{2\pi}}\mathrm{e}^{-\frac12 x_k^{2}}\right)\to 1
\end{align}
\item 当 $n\to\infty$ 时，一致地有
\begin{align}
P\!\left\{a\leqslant\frac{\mu_n-np}{\sqrt{npq}}<b\right\}\to\int_{a}^{b}\varphi(x)\,\mathrm{d}x
\end{align}
其中 $\varphi(x)=\dfrac{1}{\sqrt{2\pi}}\mathrm{e}^{-x^{2}/2}$（$-\infty<x<\infty$）。
\end{enumerate}
\end{theorem}

\begin{proof}
先证局部极限定理，我们将给出一个比上述结果更为精确的渐近式。

因 $x_k$ 只能在有限区间 $[a,b]$ 中取值，故当 $n\to\infty$ 时，
\begin{align}
k=np+x_k\sqrt{npq}\to\infty,\qquad j\equiv n-k=nq-x_k\sqrt{npq}\to\infty.
\end{align}

由 Stirling 公式：
\begin{align}
m!=\sqrt{2\pi m}\,m^{m}\mathrm{e}^{-m}\mathrm{e}^{\theta_m},\qquad\left(0<\theta_m<\frac{1}{12m}\right)
\end{align}
可知
\begin{align}
P\{\mu_n=k\}&=\frac{n!}{k!\,j!}p^{k}q^{j}\notag\\
&=\frac{\sqrt{2\pi n}\,n^{n}\mathrm{e}^{-n}}{\sqrt{2\pi k}\,k^{k}\mathrm{e}^{-k}\sqrt{2\pi j}\,j^{j}\mathrm{e}^{-j}}p^{k}q^{j}\mathrm{e}^{\theta_n-\theta_k-\theta_j}\notag\\
&=\frac{1}{\sqrt{2\pi}\sqrt{npq}}\left(\frac{np}{k}\right)^{k+\frac12}\left(\frac{nq}{n-k}\right)^{n-k+\frac12}\mathrm{e}^{\theta}
\end{align}
其中 $\theta=\theta_n-\theta_k-\theta_{n-k}$，因此
\begin{align}
|\theta|<\frac{1}{12}\left(\frac1n+\frac1k+\frac{1}{n-k}\right).
\end{align}

由前述 $k,j$ 的渐近式知
\begin{align}
\frac{k}{np}=1+x_k\sqrt{\frac{q}{np}},\qquad \frac{n-k}{nq}=1-x_k\sqrt{\frac{p}{nq}}.
\end{align}
我们将利用展开式
\begin{align}
\ln(1+x)=x-\frac{x^{2}}2+\frac{x^{3}}3-\frac{x^{4}}4+\cdots
\end{align}
来对它们进行估计。这个展开式当 $-1<x\leqslant1$ 时收敛，但只有对绝对值很小的那种 $x$ 值才收敛得快。现在，当 $n$ 充分大时，$x_k\sqrt{\dfrac{q}{np}}$ 及 $x_k\sqrt{\dfrac{p}{nq}}$ 都很小（因此，当 $p=0$ 及 $p=1$ 时不能用；此外，当 $p$ 或 $q$ 很小时，渐近展开式引起的误差也较大，这时我们已推荐用 Poisson 逼近公式）。所以
\begin{align}
&\ln\!\left(\sqrt{2\pi npq}\,P\{\mu_n=k\}\right)\notag\\
&\quad=\theta-\left(k+\frac12\right)\ln\frac{k}{np}-\left(n-k+\frac12\right)\ln\frac{n-k}{nq}\notag\\
&\quad=\theta-\left(np+x_k\sqrt{npq}+\frac12\right)\ln\!\left(1+x_k\sqrt{\frac{q}{np}}\right)\notag\\
&\qquad\quad-\left(nq-x_k\sqrt{npq}+\frac12\right)\ln\!\left(1-x_k\sqrt{\frac{p}{nq}}\right)\notag\\
&\quad=\theta-\left(np+x_k\sqrt{npq}+\frac12\right)\left(x_k\sqrt{\frac{q}{np}}-\frac{x_k^{2}q}{2np}+\frac{x_k^{3}q\sqrt{npq}}{3n^{2}p^{2}}-\frac{x_k^{4}q^{2}}{4n^{2}p^{2}}+\cdots\right)\notag\\
&\qquad\quad-\left(nq-x_k\sqrt{npq}+\frac12\right)\left(-x_k\sqrt{\frac{p}{nq}}-\frac{x_k^{2}p}{2nq}-\frac{x_k^{3}p\sqrt{npq}}{3n^{2}q^{2}}-\frac{x_k^{4}p^{2}}{4n^{2}q^{2}}+\cdots\right)\notag\\
&\quad=\theta-\frac{x_k^{2}}2+\frac{q-p}{6\sqrt{npq}}(x_k^{3}-3x_k)\notag\\
&\qquad\quad+\frac{1}{12npq}\left[3(p^{2}+q^{2})x_k^{2}-(p^{3}+q^{3})x_k^{4}\right]+o\!\left(\frac1n\right).
\end{align}
因此
\begin{align}
P\{\mu_n=k\}
&=\frac{1}{\sqrt{2\pi}}\cdot\frac{1}{\sqrt{npq}}\exp\left\{-\frac{x_k^{2}}2+\frac{(q-p)(x_k^{3}-3x_k)}{6\sqrt{npq}}+O\!\left(\frac1n\right)\right\}\notag\\
&=\frac{1}{\sqrt{2\pi}}\cdot\frac{1}{\sqrt{npq}}\mathrm{e}^{-\frac{x_k^{2}}2}\left[1+\frac{(q-p)(x_k^{3}-3x_k)}{6\sqrt{npq}}+O\!\left(\frac1n\right)\right].
\end{align}

取其第一项即得局部极限定理，因此我们已证得了局部极限定理。显然，我们得到了更精确的估计式。又知当 $p=q$ 及 $x_k^{3}-3x_k=0$ 时，近似效果尤佳。

下面转入证明积分极限定理。
\begin{align}
P\!\left\{a\leqslant\frac{\mu_n-np}{\sqrt{npq}}<b\right\}
&=P\{np+a\sqrt{npq}\leqslant\mu_n<np+b\sqrt{npq}\}\notag\\
&=\sum_{k=k_1}^{k_2}P\{\mu_n=k\},
\end{align}
上式中 $k_1$ 为不小于 $np+a\sqrt{npq}$ 的最小整数，$k_2$ 为小于 $np+b\sqrt{npq}$ 的最大整数。由局部极限定理知当 $n$ 充分大时，对任给 $\varepsilon>0$，有
\begin{align*}
P\{\mu_n=k\}=\frac{1}{\sqrt{npq}}(\varphi(x_k)+\varepsilon_k),\quad |\varepsilon_k|<\varepsilon,\qquad (k=k_1,k_1+1,\cdots,k_2).
\end{align*}
代入上式得到
\begin{align*}
P\!\left\{a\leqslant\frac{\mu_n-np}{\sqrt{npq}}<b\right\}
=\sum_{k=k_1}^{k_2}\frac{1}{\sqrt{npq}}\varphi(x_k)+\sum_{k=k_1}^{k_2}\frac{\varepsilon_k}{\sqrt{npq}}.
\end{align*}
因为有
\begin{align*}
\left|\sum_{k=k_1}^{k_2}\frac{\varepsilon_k}{\sqrt{npq}}\right|
\leqslant\frac{1}{\sqrt{npq}}(k_2-k_1+1)\varepsilon
\leqslant\frac{(b-a)\sqrt{npq}+1}{\sqrt{npq}}\varepsilon,
\end{align*}
故当 $n\to\infty$ 时，注意 $x_k$ 的增量为 $\dfrac{1}{\sqrt{npq}}$，就得到
\begin{align*}
P\!\left\{a\leqslant\frac{\mu_n-np}{\sqrt{npq}}<b\right\}\to\int_{a}^{b}\varphi(x)\,\mathrm{d}x.
\end{align*}
定理到此完全证毕。

利用 $\dfrac{1}{\sqrt{2\pi}}\int_{-\infty}^{\infty}\mathrm{e}^{-x^{2}/2}\,\mathrm{d}x=1$，不难证明在积分极限定理中，当 $a=-\infty$，$b=+\infty$ 时仍然成立。
\end{proof}

\subsection*{四、De Moivre--Laplace 极限定理的一些应用}

De Moivre--Laplace 极限定理虽然是作为二项分布的近似而推导出来的，但是它的重要性远远超出数值计算的范围。对它的各种推广形式的深入讨论将在以后几节进行，这里我们先介绍一些它的具体应用，其中有些解决了第一、二章遗留的问题。

\paragraph*{推导 Bernoulli 大数定律}
积分极限定理断言 $\mu_n$ 的分布渐近于正态分布 $N(np,npq)$，因此 Bernoulli 试验中事件 $A$ 出现的频率 $\dfrac{\mu_n}{n}$ 的分布渐近于正态分布 $N\!\left(p,\dfrac{pq}{n}\right)$。这里可以想象，当 $n\to\infty$ 时，$\dfrac{\mu_n}{n}$ 的分布会收敛于退化分布
\begin{align*}
I_p(x)=
\begin{cases}
0,&x\leqslant p,\\
1,&x>p,
\end{cases}
\end{align*}
这正是 Bernoulli 大数定律所确定的事实。下面来严格证明。

给定 $\varepsilon>0$，对任意正数 $l$，只要 $n$ 足够大，就有
\begin{align*}
l\sqrt{npq}<\varepsilon n.
\end{align*}
因此
\begin{align*}
\left\{\left|\frac{\mu_n-np}{\sqrt{npq}}\right|<l\right\}\subset\left\{\left|\frac{\mu_n-np}{n}\right|<\varepsilon\right\}.
\end{align*}
所以对大的 $n$，
\begin{align*}
P\!\left\{\left|\frac{\mu_n-np}{\sqrt{npq}}\right|<l\right\}\leqslant P\!\left\{\left|\frac{\mu_n-np}{n}\right|<\varepsilon\right\}.
\end{align*}
由积分极限定理，当 $n\to\infty$ 时，上式左边收敛于
\begin{align*}
\frac{1}{\sqrt{2\pi}}\int_{-l}^{l}\mathrm{e}^{-x^{2}/2}\,\mathrm{d}x.
\end{align*}
对于任给 $\delta>0$，可以选 $l$，使得上面积分值大于 $1-\delta$，因此对充分大的 $n$ 有
\begin{align*}
P\!\left\{\left|\frac{\mu_n}{n}-p\right|<\varepsilon\right\}>1-\delta.
\end{align*}
这就证明了 Bernoulli 大数定律，从这里也可看出，积分极限定理比大数定律更精确。

\paragraph*{用频率估计概率时的计算问题}
由积分极限定理
\begin{align}
P\!\left\{\left|\frac{\mu_n}{n}-p\right|<\varepsilon\right\}
&=P\!\left\{-\varepsilon\sqrt{\frac{n}{pq}}<\frac{\mu_n-np}{\sqrt{npq}}<\varepsilon\sqrt{\frac{n}{pq}}\right\}\notag\\
&\simeq\Phi\!\left(\varepsilon\sqrt{\frac{n}{pq}}\right)-\Phi\!\left(-\varepsilon\sqrt{\frac{n}{pq}}\right)
=2\Phi\!\left(\varepsilon\sqrt{\frac{n}{pq}}\right)-1
\end{align}
这个关系式可用来解决许多计算问题。

第一类问题是已知 $n,p,\varepsilon$，求概率 $P\!\left\{\left|\dfrac{\mu_n}{n}-p\right|<\varepsilon\right\}$；这时只要利用上式并查正态分布函数 $\Phi(x)$ 的数值表就可解决，这类问题在二项分布计算中经常会遇到。

\begin{example}[蒲丰试验]\label{example:ch05-蒲丰试验}
蒲丰试验中掷硬币 $4\,040$ 次，出正面 $2\,048$ 次，试计算当重复蒲丰试验时，正面出现的频率与概率之差的偏离程度不大于蒲丰试验中所发生的偏离的概率。
\end{example}

\begin{proof}[解]
蒲丰投币中频率 $\dfrac{2\,048}{4\,040}$ 与概率 $p=\dfrac12$ 的偏离为 $\varepsilon=\dfrac{2\,048}{4\,040}-\dfrac12=0.006\,93$，故所求概率为
\begin{align*}
\beta=P\!\left\{\left|\frac{\mu}{4\,040}-\frac12\right|\leqslant 0.006\,93\right\}
=2\Phi\!\left(0.006\,93\cdot\sqrt{\frac{4\,040}{\frac12\cdot\frac12}}\right)-1
=2\Phi(0.881\,0)-1=2\times0.810\,9-1=0.622.
\end{align*}
\end{proof}

第二类问题是要使 $\dfrac{\mu_n}{n}$ 与 $p$ 的差异不大于定数 $\varepsilon$ 的概率不小于预先给的数 $\beta$，问最少应做多少次试验？这时只需要求满足下式的最小 $n$，
\begin{align}
2\Phi\!\left(\varepsilon\sqrt{\frac{n}{pq}}\right)-1\geqslant\beta
\end{align}
这也可通过查表求得。

\begin{example}[市场占有率抽样调查]\label{example:ch05-市场占有率抽样}
某品牌往常的市场占有率为 $15\%$，今公司决定再做一次抽样调查，要求误差小于 $1\%$ 的概率达到 $95\%$，问至少要抽多少户？
\end{example}

\begin{proof}[解]
很好设计的抽样调查方案完全适合于 Bernoulli 概型。这时 $\varepsilon=0.01$，$p=0.15$，$\beta=0.95$，下面利用上式求 $n$，即
\begin{align*}
\Phi\!\left(0.01\sqrt{\frac{n}{0.15\times0.85}}\right)\geqslant\frac{\beta+1}{2}=0.975.
\end{align*}
反查标准正态分布函数的数值表得
\begin{align*}
0.01\sqrt{\frac{n}{0.15\times0.85}}\geqslant 1.96,
\end{align*}
因此
\begin{align*}
n\geqslant 4\,898.04.
\end{align*}
\end{proof}

抽样调查方案设计中，样本大小的确定至关重要，一方联系着结果的精度与可信度，另一方联系着开支预算与工作量。上例中 $p$ 有以往的信息可用，但在一般情况下只能靠估计或设计一个小型调查以提供初值。

第三类问题已知 $n$ 及 $\beta$，求 $\varepsilon$，这类问题是在进行误差估计时提出来的。解法如下：先找 $z_\beta$ 使
\begin{align*}
2\Phi(z_\beta)-1=\beta,
\end{align*}
这时
\begin{align}
\varepsilon=z_\beta\sqrt{\frac{pq}{n}}
\end{align}
即为所求。若 $p$ 不知道，则利用 $pq\leqslant\dfrac14$，有下列估计式
\begin{align}
\varepsilon\leqslant\frac{z_\beta}{2\sqrt{n}}.
\end{align}
这类估计在 Monte Carlo 方法中很有用处。

\begin{example}[抽样误差估计]\label{example:ch05-抽样误差估计}
在上例的市场占有率抽样调查中，若预算只许调查 $2\,000$ 户，可信度仍要求为 $95\%$，这时的抽样误差达到多少？
\end{example}

\begin{proof}[解]
这正是第三类计算题。这时 $n=2\,000$，$\beta=0.95$，$p=0.15$，要求 $\varepsilon$，代入公式得
\begin{align*}
\varepsilon=1.96\cdot\sqrt{\frac{0.15\times0.85}{2\,000}}=0.015\,65.
\end{align*}
假如这个误差在允许范围内，就可实行；若认为误差太大，则只有增加抽样数量，这就要求追加预算。
\end{proof}

\paragraph*{概率的置信区间估计}
由积分极限定理
\begin{align}
P\!\left\{\left|\frac{\frac{\mu_n}{n}-p}{\sqrt{\dfrac{p(1-p)}{n}}}\right|<z_\beta\right\}=\beta
\end{align}
其中 $z_\beta$ 满足 $2\Phi(z_\beta)-1=\beta$，只要 $n$ 相当大。

对 $p$ 的二次方程
\begin{align*}
\left(\frac{\mu_n}{n}-p\right)^{2}=z_\beta^{2}\frac{p(1-p)}{n}
\end{align*}
求解并略去 $\dfrac1n$ 的高阶无穷小项得：
\begin{align}
P\!\left\{\frac{\mu_n}{n}-z_\beta\sqrt{\frac{\frac{\mu_n}{n}\left(1-\frac{\mu_n}{n}\right)}{n}}<p<\frac{\mu_n}{n}+z_\beta\sqrt{\frac{\frac{\mu_n}{n}\left(1-\frac{\mu_n}{n}\right)}{n}}\right\}=\beta
\end{align}
在数理统计中称我们已在置信水平 $\beta$（一般为 $95\%$）下得到概率 $p$ 的置信区间
\begin{align}
\left(\frac{\mu_n}{n}-z_\beta\sqrt{\frac{\frac{\mu_n}{n}\left(1-\frac{\mu_n}{n}\right)}{n}},\ \frac{\mu_n}{n}+z_\beta\sqrt{\frac{\frac{\mu_n}{n}\left(1-\frac{\mu_n}{n}\right)}{n}}\right).
\end{align}

\paragraph*{局部极限定理在二项分布计算中的应用}
局部极限定理给出
\begin{align}
\binom{n}{k}p^{k}q^{n-k}\div\left(\frac{1}{\sqrt{npq}}\cdot\frac{1}{\sqrt{2\pi}}\mathrm{e}^{-\frac12\left(\frac{k-np}{\sqrt{npq}}\right)^{2}}\right)\to 1.
\end{align}
这个事实由图 5.1.1 表示出来，其中阶梯函数给出概率 $\dbinom{n}{k}p^{k}q^{n-k}$，而粗线则给出对应的正态分布密度函数曲线。

\begin{figure}[H]
\centering
\includegraphics[width=0.7\textwidth]{ch05-fig-binormal.png}
\caption{二项分布的正态逼近}
\label{figure:ch05-二项正态逼近}
\end{figure}

我们来利用局部极限定理对 $\dbinom{n}{k}p^{k}q^{n-k}$ 进行近似计算。

因为标准正态密度函数
\begin{align}
\varphi(x)=\frac{1}{\sqrt{2\pi}}\mathrm{e}^{-x^{2}/2},\quad -\infty<x<\infty
\end{align}
有专门数值表，当 $n$ 较大时，对二项分布的计算可用下列近似式：
\begin{align}
\binom{n}{k}p^{k}q^{n-k}\approx\frac{1}{\sqrt{npq}}\varphi(x_k),
\end{align}
其中 $x_k=\dfrac{k-np}{\sqrt{npq}}$。用这种方法计算当然会有误差，但实践证明（也能利用更精确的估计式作理论分析），当 $p$ 不太接近于 $0$ 或 $1$，而 $n$ 又不太小时，都能得到良好的结果。

下面就用这种办法来解决第二章 \S4 中提出的一些二项分布计算问题。

先来计算人寿保险的例 4(1)，那里要求计算 $b(40;10\,000,0.005)$，利用上式
\begin{align*}
b(40;10\,000,0.005)
&\approx\frac{1}{\sqrt{10\,000\times0.005\times0.995}}\varphi\!\left(\frac{40-10\,000\times0.005}{\sqrt{10\,000\times0.005\times0.995}}\right)\\
&=\frac{1}{7.05}\varphi(1.418)=\frac{1}{7.05}\times0.145\,6=0.020\,7,
\end{align*}
而精确值为 $0.021\,4$，误差不大。

\paragraph*{积分极限定理在二项分布计算中的应用}
积分极限定理给出
\begin{align}
P\!\left\{a\leqslant\frac{\mu_n-np}{\sqrt{npq}}<b\right\}\to\Phi(b)-\Phi(a)
\end{align}
其中 $\Phi(x)=\dfrac{1}{\sqrt{2\pi}}\int_{-\infty}^{x}\mathrm{e}^{-t^{2}/2}\,\mathrm{d}t$ 是标准正态分布函数，利用它可以计算形如 $P\{k_1\leqslant\mu_n\leqslant k_2\}$ 的概率。

因为当 $n$ 充分大时，$b(k;n,p)$ 即使对于最可能成功次数，其数值也很小，这时讨论成功次数等于某数 $k$ 的概率 $P\{\mu_n=k\}$ 便没有太大意思；相反，计算 $\mu_n$ 的值落在某一区间 $[k_1,k_2]$ 的概率倒十分重要，因此积分极限定理在二项分布近似计算中更为有用。

我们有
\begin{align}
P\{k_1\leqslant\mu_n\leqslant k_2\}
&=P\!\left\{\frac{k_1-np}{\sqrt{npq}}\leqslant\frac{\mu_n-np}{\sqrt{npq}}\leqslant\frac{k_2-np}{\sqrt{npq}}\right\}\notag\\
&\approx\Phi\!\left(\frac{k_2-np}{\sqrt{npq}}\right)-\Phi\!\left(\frac{k_1-np}{\sqrt{npq}}\right).
\end{align}

细察图 5.1.1 可以看出，在用连续型分布正态分布逼近离散型分布二项分布时存在着偏差。对应于离散值 $7,8,9,10$ 的矩形面积，相应的该是 $6.5<x<10.5$ 上的曲边梯形面积，因此下面的修正公式，常能得到更好的近似效果：
\begin{align}
P\{k_1\leqslant\mu_n\leqslant k_2\}
\approx\Phi\!\left(\frac{k_2-np+0.5}{\sqrt{npq}}\right)-\Phi\!\left(\frac{k_1-np-0.5}{\sqrt{npq}}\right).
\end{align}

当 $p$ 不太接近 $0$ 或 $1$，而 $n$ 又不太小时，用这个近似式能得到良好的结果。

特别当 $k_1=k_2$ 时，上式化为局部极限定理的近似式，与局部极限定理的结果完全统一。

例如，我们来完成第二章 \S4 中的一些计算。人寿保险的例 4(2) 要计算 $P\{\mu\leqslant70\}$，利用修正公式
\begin{align*}
P\{\mu\leqslant70\}
&\approx\Phi\!\left(\frac{70-10\,000\times0.005+0.5}{\sqrt{10\,000\times0.005\times0.995}}\right)-\Phi\!\left(\frac{-10\,000\times0.005-0.5}{\sqrt{10\,000\times0.005\times0.995}}\right)\\
&=\Phi(2.91)-\Phi(-7.16)=\Phi(2.91)-[1-\Phi(7.16)]=0.998.
\end{align*}
这里 $\Phi(7.16)$ 取为 $1$，此外计算中还用到 $\Phi(-x)=1-\Phi(x)$，这由 $\varphi(x)$ 的对称性即得。

可见，保险公司有很大把握假设死亡人数不大于 $70$ 人，并可据此作各种估算。当然 $70$ 是举例性的，读者不妨试着计算 $P\{\mu\leqslant100\}$。

第二章 \S4 的例 5，即机票超售问题，作为贯穿本书的一个案例，是很现实的一个问题。为使问题能在允许的篇幅内有答案，已被数学模型化，假定各个旅客登机的独立性，归为 Bernoulli 概型，并把问题的要求加以简化与明确，即对不同的超售额 $m$ 计算发生麻烦（即有旅客被拒登机）的概率。这是在用数学解决实际问题时必须经过的一步，也是没有标准答案的一步，因此不同理论修养和实践经验的应用数学家将采用大相径庭的数学模型来处理同一个实际问题。

我们的方案是对适当的 $m$ 计算如下概率：
\begin{align}
P=P\{\mu>200\}=\sum_{k>200}b(k;200+m,0.95).
\end{align}
有了积分极限定理以及近似计算公式，任务不难完成。例如当 $m=7$ 时，
\begin{align*}
P=\sum_{k=201}^{207}b(k;207,0.95)
\approx\Phi\!\left(\frac{207-207\times0.95+0.5}{\sqrt{207\times0.95\times0.05}}\right)
-\Phi\!\left(\frac{201-207\times0.95-0.5}{\sqrt{207\times0.95\times0.05}}\right)
=\Phi(3.460)-\Phi(1.228)=0.9997-0.8903=0.1094.
\end{align*}

相应于 $m=1,2,3,4,5,6,7$ 的 $P$ 值已列于第一章 \S1 中，供主管部门使用。实际问题的最终解决大都要求应用数学家与实际部门反复磋商。显然上式也可用 Poisson 逼近计算。

现在我们能够解决第二章 \S4 例 6 中提出的车间用电问题了。该问题是要求 $r$，使
\begin{align}
P\{\mu\leqslant r\}=\sum_{k=0}^{r}\binom{200}{k}(0.6)^{k}(0.4)^{200-k}\geqslant 0.999.
\end{align}
我们可以利用积分极限定理计算这个概率。
\begin{align*}
\sum_{k=0}^{r}\binom{200}{k}(0.6)^{k}(0.4)^{200-k}
&\approx\Phi\!\left(\frac{r-200\times0.6+0.5}{\sqrt{200\times0.6\times0.4}}\right)-\Phi\!\left(\frac{-200\times0.6-0.5}{\sqrt{200\times0.6\times0.4}}\right)\\
&=\Phi\!\left(\frac{r-119.5}{\sqrt{48}}\right)-\Phi(-17.39)\approx\Phi\!\left(\frac{r-119.5}{\sqrt{48}}\right)\geqslant 0.999.
\end{align*}
查表得
\begin{align*}
\frac{r-119.5}{\sqrt{48}}=3.09,
\end{align*}
所以
\begin{align*}
r=141.
\end{align*}

这个结果表明 $P\{\mu\leqslant141\}\geqslant0.999$，所以我们若供电 $141$ 千瓦，那么由于供电不足而影响生产的可能性小于 $0.001$，相当于在 $8$ 小时工作中有半分钟受影响，这在一般工厂中是允许的。当然不同的生产单位，可能提出不同的要求，那么我们可以改变右端的概率值，但是方法还是同样的。

再来计算作为分子运动模型的例 7，其中 $p=\dfrac12$，$n=5.4\times10^{22}$。
\begin{align*}
P\{|\mu-np|>2.7\times10^{12}\}
&=P\!\left\{\frac{|\mu-np|}{\sqrt{npq}}>\frac{2.7\times10^{12}}{\sqrt{5.4\times10^{22}\times\frac14}}\right\}\\
&=P\!\left\{\frac{|\mu-np|}{\sqrt{npq}}>23.2\right\}\simeq\frac{2}{\sqrt{2\pi}}\int_{23.2}^{\infty}\mathrm{e}^{-t^{2}/2}\,\mathrm{d}t.
\end{align*}
这个数值非常小，从 $\Phi(x)$ 的数值表中不能找到，但我们可以用下面方法对它进行估计。

由于对 $z>0$ 有
\begin{align}
\int_{z}^{\infty}\mathrm{e}^{-t^{2}/2}\,\mathrm{d}t<\frac1z\int_{z}^{\infty}t\mathrm{e}^{-t^{2}/2}\,\mathrm{d}t=\frac1z\mathrm{e}^{-z^{2}/2}.
\end{align}
利用它可以得到
\begin{align*}
\frac{2}{\sqrt{2\pi}}\int_{23.2}^{\infty}\mathrm{e}^{-t^{2}/2}\,\mathrm{d}t
<\frac{2}{\sqrt{2\pi}}\cdot\frac{1}{23.2}\mathrm{e}^{-269.12}<10^{-100}.
\end{align*}
这概率非常非常之小。

这种计算为分子运动论提供佐证：虽然每个气体分子的运动轨道、速度、方向都是随机的，但因为分子数量十分巨大，从宏观上看，作为平均特征的压力、温度等却是十分稳定的。

现在也能对第一章 \S1 中的高尔顿板这个试验作理论解释了。读者可能早已想到，近似小球高度的曲线是正态分布密度函数曲线。这完全正确。事实上，高尔顿板可以看作是 Bernoulli 试验的一个实验模型。如果我们把小球碰到钉子看作是一次试验，而把从右边落下算是成功，当然从左边落下就算是失败，这时就有了一次 $p=\dfrac12$ 的 Bernoulli 试验。小球从顶端到底层共需要经过 $n$ 排钉子，这就相当于一个 $n$ 次 Bernoulli 试验，剩下的只是要说明为什么高度曲线会是正态分布密度函数曲线，这个问题留给读者思考。

\end{document}
